
Out-of-memory (OOM) failures during neural network training waste computational resources. A researcher training ResNet-34 with batch size 64 faces trial-and-error batch size tuning, with GPU-hours wasted on failed experiments. Without accurate pre-training memory prediction, OOM failures are discovered after hours of setup rather than predicted in seconds.

Existing lightweight memory profiling methods have made progress. VeritasEst (2025) demonstrated that 2-iteration profiling suffices for optimizer state stabilization, achieving 5.46\% median error through CPU-based allocator simulation. Kang et al. (2025) refined memory estimation through factorization into model parameters, gradients, optimizer states, and activations, achieving 8.7\% mean error on multimodal models.

These methods miss a critical timing window. VeritasEst's 2-iteration protocol samples memory before \texttt{optimizer.step()}, capturing gradients but missing the moment when Adam allocates momentum buffers. Optimizer workspace allocations occur during \texttt{optimizer.step()}, not during backward propagation. For Adam and AdamW optimizers, which maintain first-order moment (m\_t) and second-order moment (v\_t) buffers, this represents approximately $2\times$ parameter memory allocated when prior methods stop measuring.

Our key insight emerged from empirical observation during experiments with ResNet-18 and Adam optimizer. Memory increased from 130MB (after backward pass) to 280MB (after \texttt{optimizer.step()})---a 150MB allocation representing Adam's m\_t and v\_t workspace buffers. This allocation occurred in the gap between prior measurement points. Sampling after \texttt{optimizer.step()} completes captures this workspace allocation; sampling before misses it entirely.

We designed a post-optimizer 3-iteration sampling protocol. The first iteration runs a forward-only pass, establishing baseline memory for model parameters and forward activations. The second measurement occurs immediately after the first backward pass and \texttt{optimizer.step()}, capturing peak memory inclusive of optimizer workspace. The prediction combines measurements via max operation: \texttt{max(iter1\textbackslash \_forward, post\textbackslash \_optim\textbackslash \_peak)}.

Validating this protocol on ResNet-18 and ResNet-34 with Adam and SGD optimizers demonstrates its effectiveness. Across 4 configurations using synthetic data (generated via \texttt{torch.randn()} to match CIFAR-10 dimensions), post-optimizer sampling achieves 2.6\% median relative error---52\% lower than the 5.46\% baseline. All tested configurations remain below 7\% error, with 95th percentile error at 6.1\%. The protocol exhibits optimizer-specific accuracy: Adam achieves 0.0-0.62\% error, while SGD achieves 4.6-6.4\% error. This $17\times$ accuracy difference reveals that profiling effectiveness depends on optimizer type.

This work makes three primary contributions. First, we introduce the post-optimizer sampling protocol, which achieves 52\% error reduction versus state-of-the-art by timing measurements to capture Adam workspace allocations. Second, we provide empirical demonstration that memory profiling accuracy varies substantially by optimizer type, with Adam showing $17\times$ better accuracy than SGD for identical architectures. Third, we establish optimizer workspace---not forward activations---as the critical memory component for accurate profiling, reframing the design space for future memory prediction systems.

This work is scoped to CNN architectures (ResNet-18/34) with fixed-length inputs, validated using synthetic data. Validation used Adam and SGD optimizers on CIFAR-10-scale dimensions ($32\times 32\times 3)$. Transformer architectures, variable-length sequence inputs, and real dataset validation remain future work. The reduced validation scale (4 configurations versus the planned 48) prioritizes mechanism validation over statistical power.

The following section surveys related work in lightweight memory profiling and positions our timing optimization contribution relative to existing approaches.

